\documentclass[11pt,a4paper]{article}
\usepackage[T1]{fontenc}
\usepackage{lmodern,amsmath,amssymb,amsthm,mathtools,booktabs,microtype}
\usepackage[margin=26mm]{geometry}
\usepackage[colorlinks=true,linkcolor=blue,urlcolor=blue]{hyperref}
\setlength{\emergencystretch}{2em}
\newtheorem{theorem}{Theorem}[section]
\newtheorem{lemma}[theorem]{Lemma}
\newtheorem{corollary}[theorem]{Corollary}
\newcommand{\Q}{\mathbb Q}
\title{A three-coordinate quadratic chart for an elliptic point and its translate}
\author{}\date{15 November 2026}
\begin{document}\maketitle
\begin{abstract}
Two quadratic equations in three field coordinates describe a point on
$y^2=x^3-Dx$ and its translate by a fixed point. The chart covers exactly
the locus on which both points are finite with nonzero abscissas. Three
linear functions are units on the entire chart, with explicit polynomial
inverses. The proof includes tangent intersections, two-torsion points
with nonzero abscissa, and every other exceptional secant configuration.
No rank or arithmetic realization hypothesis is needed.
\end{abstract}

\section{The equations and the full domain}
Let $k$ be a field of characteristic different from two. Fix
$a,b,D\in k$ such that
\[
 abD\ne0,\qquad b^2=a^3-Da,
 \qquad c=b/a,\qquad \delta=D/a=a-c^2.
\]
Write
\[
 E:y^2=x^3-Dx,\qquad P=(a,b),\qquad T_0=(0,0).
\]
The hypothesis $D\ne0$ makes $E$ nonsingular. Define the affine curve
$C\subset\mathbb A_k^3$, with coordinates $S,T,v$, by
\begin{equation}\label{eq:chart}
 \begin{split}
 Q_1&=v(v+T-S)+\delta=0,\\
 Q_2&=ST+c(T-S)+2cv+2\delta=0.
 \end{split}
\end{equation}
Set
\begin{equation}\label{eq:linear}
 U=v-S-c,\qquad W=v+T-c,\qquad
 u=U/a,\quad w=W/a,\quad\mu=v-c.
\end{equation}

\begin{theorem}\label{thm:chart}
The curve $C$ is isomorphic over $k$ to
\[
                       E\setminus\{O,T_0,P,P+T_0\}.
\]
At every point of $C$ the functions $u,w,v$ are nonzero and satisfy
\begin{equation}\label{eq:units}
 uw=1/a,\qquad
 v^{-1}=-\frac{v+T-S}{\delta}
       =\frac{v-a(u+w)-2c}{\delta}.
\end{equation}
The corresponding point and translate are
\begin{equation}\label{eq:recovery}
 H=\left(\frac vu,S\frac vu\right),\qquad
 H-P=\left(\frac vw,T\frac vw\right).
\end{equation}
In particular
\[
 \frac1{x(H)}=\frac uv,
 \qquad\frac1{x(H-P)}=\frac wv,
 \qquad S=\frac{y(H)}{x(H)},\quad T=\frac{y(H-P)}{x(H-P)}.
\]
No exclusion of $H=-P$, of a nonzero-abscissa two-torsion point,
or of a repeated line intersection is required beyond the displayed
four-point complement.
\end{theorem}

\section{Algebraic recovery and the group law}
The elementary identity
\begin{equation}\label{eq:uw}
                            UW-a=Q_1-Q_2
\end{equation}
gives $uw=1/a$. Equation $Q_1=0$ and $\delta\ne0$ give $v\ne0$
and its first inverse in~\eqref{eq:units}. The second follows from
$a(u+w)=2v+T-S-2c$. These statements hold in the coordinate ring,
not only on a dense open subset.

Put $x=v/u$ and $x'=v/w$. All divisions are now justified. From
$uw=1/a$ and $Q_1=0$ one obtains
\begin{equation}\label{eq:sum-product}
 x+x'=\mu^2-a,\qquad xx'=av^2.
\end{equation}
Indeed $x+x'=av(u+w)=v(U+W)$, and
$v(U+W)-(\mu^2-a)=Q_1$. The second equality follows directly from
$uw=1/a$.

Consider the nonvertical line
\[
                             y=\mu x-av.
\]
At $x=a$ its ordinate is $-ac=-b$, so it passes through $-P$.
The identity
\begin{equation}\label{eq:cubic}
 Z^3-DZ-(\mu Z-av)^2
   =(Z-a)\bigl(Z^2-(\mu^2-a)Z+av^2\bigr)
\end{equation}
uses only $D=a(a-c^2)$ and $\mu=v-c$.
Equations~\eqref{eq:sum-product} identify its other two roots as
$x,x'$, counted with multiplicity. The definitions give
\[
 Sx=\mu x-av,\qquad -Tx'=\mu x'-av.
\]
Thus the three intersections of the line with $E$ are
$H=(x,Sx)$, $-P$, and $-H'= (x',-Tx')$, with intersection
multiplicity. The chord-and-tangent group law yields
$H-P-H'=O$, proving $H'=H-P$.

Both recovered points are finite, and their abscissas are nonzero
because $u,w,v$ are units. Consequently $H$ lies in the asserted open
set. In particular $H=P$ cannot arise: at abscissa $a$ the line's
ordinate is $-b$, which differs from $b$ because $2b\ne0$.

The multiplicity statement treats all exceptional configurations.
If $x=a$, the point $H$ is $-P$ and the line is tangent there when
the root is repeated. If $x=x'$, the line is tangent at $H=-H'$.
A triple root is the corresponding flex configuration, and the same
group-law identity holds. At an allowed two-torsion point the ratio
$S$ or $T$ is zero; no step divides by either ratio or by a variable
ordinate. This proves the forward direction of
Theorem~\ref{thm:chart} without an additional exceptional constant.

\section{Converse and regularity}
Take $H$ in the four-point complement, and put $H'=H-P$.
The line through $H$ and $-P$ is nonvertical. Indeed a vertical line
would have $x(H)=a$, and then $H=P$ or $-P$; the first is excluded,
while in the second case the tangent at $-P$ is nonvertical since
$b\ne0$. Let its slope be $\mu$ and put $v=\mu+c$. Its equation is
$y=\mu x-av$.

By the group law, its third intersection is $-H'$. If
$x=x(H)$ and $x'=x(H')$, comparison with~\eqref{eq:cubic} gives
\[
                         xx'=av^2,\qquad x+x'=\mu^2-a.
\]
The abscissas and $a$ are nonzero, so $v\ne0$. Define
\[
 S=y(H)/x,\qquad T=y(H')/x',\qquad u=v/x,\qquad w=v/x'.
\]
The line equation gives $S=v-c-au$ and $T=aw-v+c$, and the product
identity gives $uw=1/a$. The sum identity now implies $Q_1=0$.
Identity~\eqref{eq:uw} then implies $Q_2=0$.

These formulas are mutually inverse to~\eqref{eq:recovery}. They
also define regular maps on the entire indicated loci. The reciprocal
abscissas and the two ratios are regular on the four-point complement.
For the line slope, the usual expression
\[
                         \mu=\frac{y(H)+b}{x(H)-a}
\]
is regular wherever $x(H)\ne a$. Near the only remaining point
$H=-P$, the curve equation gives the equal expression
\[
             \mu=\frac{x(H)^2+a x(H)+a^2-D}{y(H)-b}.
\]
Its denominator equals $-2b$ at $-P$ and is nonzero there. Thus the
tangent value is regular, not a separately chosen branch. On $C$,
the polynomial unit inverses in~\eqref{eq:units} and
$u^{-1}=aw$, $w^{-1}=au$ make every recovery expression regular.
This proves the isomorphism assertion.

\section{An equivalent chart displaying the units}
The invertible affine change~\eqref{eq:linear} gives an equivalent
description using $u,w,v$ as the three coordinates:
\begin{equation}\label{eq:unit-chart}
              auw=1,\qquad
              v\bigl(v-a(u+w)-2c\bigr)=\delta.
\end{equation}
The ratios are the linear expressions
$S=v-c-au$ and $T=aw-v+c$. In particular all three inverses are
linear functions of the chart coordinates. The two reciprocal
abscissas are the quadratic functions
\[
 \frac1{x(H)}=\frac{u\bigl(v-a(u+w)-2c\bigr)}{\delta},
 \qquad
 \frac1{x(H-P)}=\frac{w\bigl(v-a(u+w)-2c\bigr)}{\delta}.
\]
Their squares are generally of degree four in these coordinates.
An application which requires stored reciprocal squares must account
for that degree separately; Theorem~\ref{thm:chart} specifies the
reference point chart and its exact domain.

For the elliptic data used in the defining-polynomial constructions,
\[
 k=\Q(\rho),\quad\rho^2=2,\quad D=11-6\rho,
 \quad a=\frac{57}{8}-\frac72\rho,
 \quad b=-\frac{165}{8}+\frac{417}{32}\rho,
\]
the constants are
\[
 c=-\frac{87}{41}+\frac{129}{164}\rho,
 \qquad \delta=\frac{2328-272\rho}{1681}.
\]
Thus~\eqref{eq:chart} is exactly four rational quadratic equations
in six rational coordinates. Its coefficients have a finite explicit
specification without solving any variable linear equation.

\section{Polynomial certificates and exact checks}
The following identities give short independent certificates for the
recovered curve equations. With the notation above, put
\[
 J_U=av-(\mu^2-a)U+vU^2,\qquad
 J_W=av-(\mu^2-a)W+vW^2.
\]
Then
\begin{align*}
 J_U&=-(S+c)Q_1+vQ_2,\\
 J_W&=(T-c)Q_1+vQ_2,\\
 UvS^2+\delta U^2-av^2&=(S+c)^2Q_1-v(S+c)Q_2,\\
 WvT^2+\delta W^2-av^2&=(T-c)^2Q_1+v(T-c)Q_2.
\end{align*}
After multiplication by the relevant already nonzero denominators,
the last two identities are precisely the equations of $H$ and $H'$.
Together with~\eqref{eq:uw} and~\eqref{eq:cubic}, they are verified
by ordinary coefficient arithmetic in $\Q[a,c,S,T,v,Z]$, with
$D=a(a-c^2)$.

The standard-library checker
\nolinkurl{computation/elliptic_charts/three_coordinate_v40/check.py}
verifies these universal identities independently, authenticates the
specialized quadratic-field constants, and checks the maps in both
directions on exact project points, including $H=-P$, allowed
two-torsion points, and $2H=P$. It also exhausts the complete chart
over the stated small finite fields for every allowed fixed $P,D$,
comparing with an independent elliptic group implementation. These
finite checks supplement the all-field proof above; they are not
an assertion about any prime-pattern or height premise.
\end{document}
